\chapter{Design}

The design phase translates functional requirements into precise architectural blueprints, interaction models, flowcharts, data schemas, and interface structures.

\section{System Architecture}
LegalAssist uses a modular three-tier web architecture consisting of the Presentation Layer, Application Layer, and Data Layer.

\begin{figure}[h!]
    \centering
    \begin{tikzpicture}[node distance=2.2cm, auto]
        \node (pres) [process, fill=blue!15, minimum width=12cm, minimum height=1.2cm] {\textbf{Presentation Layer (React + Vite + Bootstrap)}\\Pages, Components, Contexts, Navigation};
        \node (app) [process, fill=purple!15, below of=pres, minimum width=12cm, minimum height=1.2cm] {\textbf{Application Layer (Node.js + Express REST APIs)}\\Auth Middleware, Availability Logic, Slot Locking, Case Pipeline};
        \node (data) [process, fill=green!15, below of=app, minimum width=12cm, minimum height=1.2cm] {\textbf{Data Layer (Hybrid Persistence)}\\Persistent JSON Engine / MongoDB Atlas Cloud Synchronization};
        
        \draw [arrow, <->] (pres) -- node [right] {HTTP / REST / JSON} (app);
        \draw [arrow, <->] (app) -- node [right] {Data I/O / Sync} (data);
    \end{tikzpicture}
    \caption{System Architecture of LegalAssist}
    \label{fig:system_architecture}
\end{figure}

\begin{enumerate}
    \item \textbf{Presentation Layer:} Built with React, Vite, and Bootstrap, rendering responsive dashboards tailored to Client, Advocate, and Admin roles.
    \item \textbf{Application Layer:} Built with Express.js, providing secure RESTful API endpoints for authentication, advocate approval, availability management, slot validation, case stage updates, AI queries, and SOS triggers.
    \item \textbf{Data Layer:} Utilizes a hybrid data access layer that operates on a local persistent JSON engine with optional automatic sync to MongoDB Atlas collections.
\end{enumerate}

\section{Use Case Design}
The system defines interactions across three primary actors: Client (User), Advocate, and Administrator.

\begin{figure}[h!]
    \centering
    \begin{tikzpicture}
        % Actors
        \node (client) [actor] {Client};
        \node (advocate) [actor, below=3cm of client] {Advocate};
        \node (admin) [actor, below=3cm of advocate] {Admin};
        
        % Use Cases
        \node (uc1) [usecase, right=3cm of client] {Register / Login};
        \node (uc2) [usecase, below=0.5cm of uc1] {Browse Advocates};
        \node (uc3) [usecase, below=0.5cm of uc2] {Book Slot};
        \node (uc4) [usecase, below=0.5cm of uc3] {Track Case \& Docs};
        \node (uc5) [usecase, below=0.5cm of uc4] {Set Availability};
        \node (uc6) [usecase, below=0.5cm of uc5] {Approve / Reject Slot};
        \node (uc7) [usecase, below=0.5cm of uc6] {Verify Advocates};
        
        % Connections Client
        \draw (client) -- (uc1);
        \draw (client) -- (uc2);
        \draw (client) -- (uc3);
        \draw (client) -- (uc4);
        
        % Connections Advocate
        \draw (advocate) -- (uc1);
        \draw (advocate) -- (uc5);
        \draw (advocate) -- (uc6);
        \draw (advocate) -- (uc4);
        
        % Connections Admin
        \draw (admin) -- (uc1);
        \draw (admin) -- (uc7);
    \end{tikzpicture}
    \caption{Use Case Diagram of LegalAssist}
    \label{fig:use_case}
\end{figure}

\section{Data Flow Diagram (DFD)}
DFDs model the flow of information between external entities, system processes, and persistent data stores.

\subsection{Level 0 Data Flow Diagram}
\begin{figure}[h!]
    \centering
    \begin{tikzpicture}[node distance=3cm]
        \node (u) [process, fill=blue!10] {Client / Advocate / Admin};
        \node (sys) [process, circle, fill=orange!20, right of=u, minimum size=2.5cm] {\textbf{0.0 LegalAssist System}};
        \node (db) [database, right of=sys] {Data Store};
        
        \draw [arrow] (u) -- node [above] {Requests / Input} (sys);
        \draw [arrow] (sys) -- node [below] {Responses / Data} (u);
        \draw [arrow] (sys) -- node [above] {Write Data} (db);
        \draw [arrow] (db) -- node [below] {Fetch Data} (sys);
    \end{tikzpicture}
    \caption{Level 0 Data Flow Diagram}
    \label{fig:dfd0}
\end{figure}

\subsection{Level 1 Data Flow Diagram}
\begin{figure}[h!]
    \centering
    \begin{tikzpicture}[node distance=2cm]
        \node (p1) [process] {1.0 Auth \& Verification};
        \node (p2) [process, right=1.5cm of p1] {2.0 Availability \& Slots};
        \node (p3) [process, below=1.5cm of p1] {3.0 Booking \& Appointments};
        \node (p4) [process, right=1.5cm of p3] {4.0 Case \& Document Pipeline};
        
        \node (d1) [database, right=1.5cm of p2] {Users DB};
        \node (d2) [database, right=1.5cm of p4] {Appointments \& Cases DB};
        
        \draw [arrow] (p1) -- (d1);
        \draw [arrow] (p2) -- (d1);
        \draw [arrow] (p3) -- (d2);
        \draw [arrow] (p4) -- (d2);
    \end{tikzpicture}
    \caption{Level 1 Data Flow Diagram}
    \label{fig:dfd1}
\end{figure}

\section{Activity Design}
Activity flow diagrams describe control execution during critical processes.

\subsection{User Authentication Activity}
The authentication activity handles login verification, role validation, and status checks (e.g., verifying that an advocate's account has been approved by an administrator before granting dashboard access).

\begin{figure}[h!]
    \centering
    \begin{tikzpicture}[node distance=1.3cm]
        \node (start) [process, fill=red!20] {Start};
        \node (input) [process, fill=blue!20, below of=start] {Enter Email, Password \& Role};
        \node (check) [decision, below of=input] {Credentials Valid?};
        \node (role) [decision, below of=check, node distance=2.2cm] {Is Advocate?};
        \node (status) [decision, below of=role, node distance=2.2cm] {Status Verified?};
        \node (auth) [process, fill=green!20, below of=status, node distance=2cm] {Grant Access \& Token};
        \node (err) [process, fill=red!20, right=2cm of check] {Show Error Message};
        \node (stop) [process, fill=red!20, below of=auth] {End};
        
        \draw [arrow] (start) -- (input);
        \draw [arrow] (input) -- (check);
        \draw [arrow] (check) -- node[above]{No} (err);
        \draw [arrow] (check) -- node[right]{Yes} (role);
        \draw [arrow] (role) -- node[right]{Yes} (status);
        \draw [arrow] (role) -- node[above]{No} (auth);
        \draw [arrow] (status) -- node[right]{Yes} (auth);
        \draw [arrow] (status) -| node[above]{No} (err);
        \draw [arrow] (auth) -- (stop);
    \end{tikzpicture}
    \caption{User Authentication Activity Diagram}
    \label{fig:activity_auth}
\end{figure}

\section{Database Design}
The database engine maintains structured entities for Users, Advocates, Appointments, Cases, Availability, SOS Alerts, and Reviews.

\begin{figure}[h!]
    \centering
    \begin{tikzpicture}[node distance=1cm]
        \node (root) [process, fill=blue!10] {\textbf{legalassist\_db}};
        \node (users) [process, below left=1.2cm and 2cm of root] {\textbf{users}\\id, name, email, role, status, barId, fees, availability};
        \node (app) [process, below left=1.2cm and -0.5cm of root] {\textbf{appointments}\\id, advocateId, userId, date, time, status, reason};
        \node (cases) [process, below right=1.2cm and -0.5cm of root] {\textbf{cases}\\id, clientId, advocateId, stage, documents};
        \node (sos) [process, below right=1.2cm and 2cm of root] {\textbf{sosAlerts / reviews}\\id, userId, message, rating, comment};
        
        \draw [arrow] (root) -- (users);
        \draw [arrow] (root) -- (app);
        \draw [arrow] (root) -- (cases);
        \draw [arrow] (root) -- (sos);
    \end{tikzpicture}
    \caption{Database Conceptual Schema}
    \label{fig:db_schema}
\end{figure}

\section{Security \& Authorization Design}
Security is enforced using role-based access control (RBAC), strict server-side request validation, input sanitization, and atomic slot reservation guards preventing duplicate concurrent bookings.

\clearpage
